\documentclass{article}

% NeurIPS 2026 style was not yet publicly distributed at the time of writing;
% this uses the most recent official NeurIPS style file (2025) in preprint mode.
\usepackage[preprint]{neurips_2025}

\usepackage[utf8]{inputenc}
\usepackage[T1]{fontenc}
\usepackage{hyperref}
\usepackage{url}
\usepackage{booktabs}
\usepackage{amsfonts}
\usepackage{amsmath}
\usepackage{microtype}
\usepackage{xcolor}

% Every number in Section 6 is generated from the committed run artifacts by
% paper/gen_macros.py; none are typed by hand.
% Auto-generated by gen_macros.py from benchmark/results/*.json — do not edit by hand.
\newcommand{\NPrompts}{541}
\newcommand{\NPromptsPaired}{541}
\newcommand{\NInstrInstances}{834}
\newcommand{\ControlPromptStrict}{87.1}
\newcommand{\TreatmentPromptStrict}{93.9}
\newcommand{\DeltaPromptStrict}{6.8}
\newcommand{\DeltaCILow}{4.6}
\newcommand{\DeltaCIHigh}{9.1}
\newcommand{\McNemarP}{7.8 \times 10^{-10}}
\newcommand{\DiscordantTreatOnly}{39}
\newcommand{\DiscordantControlOnly}{2}
\newcommand{\ControlPromptLoose}{90.4}
\newcommand{\TreatmentPromptLoose}{95.4}
\newcommand{\DeltaPromptLoose}{5.0}
\newcommand{\ControlInstrStrict}{91.0}
\newcommand{\TreatmentInstrStrict}{96.0}
\newcommand{\GateInitialFail}{61}
\newcommand{\GateFinalFail}{22}
\newcommand{\GateSelectedRevision}{61}
\newcommand{\ExtraCalls}{90}
\newcommand{\PromptsWithRevision}{61}
\newcommand{\MeanRevisions}{0.17}
\newcommand{\MeanCharsControl}{1537}
\newcommand{\MeanCharsTreatment}{1546}
\newcommand{\ExactN}{451}
\newcommand{\ExactPct}{54.1}
\newcommand{\ExactControl}{94.2}
\newcommand{\ExactTreat}{98.4}
\newcommand{\ExactDelta}{4.2}
\newcommand{\ApproxN}{288}
\newcommand{\ApproxPct}{34.5}
\newcommand{\ApproxControl}{89.6}
\newcommand{\ApproxTreat}{97.6}
\newcommand{\ApproxDelta}{8.0}
\newcommand{\UncoveredN}{95}
\newcommand{\UncoveredPct}{11.4}
\newcommand{\UncoveredControl}{80.0}
\newcommand{\UncoveredTreat}{80.0}
\newcommand{\UncoveredDelta}{0.0}
\newcommand{\AllN}{834}
\newcommand{\AllPct}{100.0}
\newcommand{\AllControl}{91.0}
\newcommand{\AllTreat}{96.0}
\newcommand{\AllDelta}{5.0}


\title{Charter: Declarative Deliverable Contracts for Governing Long-Horizon Agent Runs}

\author{%
  The \texttt{agent-charter} contributors\\
  \texttt{github.com/zhenthebuilder/agent-charter}\\
  \textit{This paper was produced and evaluated during an autonomous research run.}
}

\begin{document}

\maketitle

\begin{abstract}
  Long-horizon agent tasks fail quietly: the agent finishes, sounds confident,
  and delivers artifacts that are missing sections, contain leftover scaffolding,
  or silently drop requirements stated hours earlier in the prompt. We argue that
  the deliverables of such runs should be governed the way software releases are
  governed---by a small, declarative, machine-checkable contract---and we present
  \textsc{Charter}, an open-source, dependency-free tool that turns a list of
  required artifacts and acceptance checks into (i) instructions for the agent,
  (ii) a deterministic gate with a CI-style exit code, and (iii) an audit report
  with file hashes. To quantify the value of gating generation on a deliverable
  contract, we run a paired A/B on IFEval, a public benchmark of \NPrompts{}
  prompts with verifiable output constraints: single-shot generation versus the
  same model revised under a \textsc{Charter} gate compiled automatically from
  each prompt's constraints. The gate raises prompt-level strict accuracy from
  \ControlPromptStrict\% to \TreatmentPromptStrict\%
  (+\DeltaPromptStrict{} points, McNemar $p = \McNemarP$) at a cost of
  \MeanRevisions{} additional model calls per prompt on average, and the gain is
  concentrated almost entirely on constraints the contract can express---%
  constraint types outside the contract's coverage are essentially unchanged.
  This both demonstrates the practical value of contract-gated generation and
  delineates its boundary: a gate can only fix what it can check. Code, contracts,
  raw logs, and scoring scripts are released for exact reproduction.
\end{abstract}

\section{Introduction}

Agents built on large language models are now routinely given tasks that run for
hours: produce a research report with its underlying data, implement a feature
with tests and documentation, analyze a dataset and write up the findings. The
economics of these systems depend on the human \emph{not} supervising every step.
Yet the dominant failure mode of long-horizon runs is not a crash that pages
someone---it is a confident ``done'' attached to deliverables that do not match
what was asked: a missing section, a \texttt{TODO} left on page four, a data file
that was never written, tests that never ran, instructions stated at hour zero
that fell out of effective context by hour three.

Today the main control mechanism is post-hoc human review of everything the agent
produced, which does not scale with agent throughput, or ad-hoc prompt
engineering (``make sure the report has a limitations section''), which is not
enforced by anything. LLM-as-a-judge pipelines~\citep{zheng2023judging} help with
fuzzy quality criteria but are stochastic, add cost to every iteration, and are
hard to audit as a release gate. What is missing is the analogue of what made
software delivery governable: a small, reviewable specification of ``done'' that
a machine can check on every iteration, cheaply and deterministically.

We make that analogy concrete. \textsc{Charter} is a deliverable contract for
agent runs: a declarative TOML/JSON file that lists the artifacts a run must
produce and the acceptance checks each must pass---existence, size and length
bounds, required sections, required or forbidden patterns, valid JSON with
required keys, absence of placeholder text, and arbitrary commands (test suites,
linters, schema validators). The contract is used three ways by one tool:
\texttt{charter explain} renders it as instructions injected into the agent's
prompt; \texttt{charter check} evaluates the workspace and returns a CI-style
exit code, with a structured failure report designed to be fed back to the model;
\texttt{charter report} produces an audit trail (per-check verdicts, file hashes,
timestamps) that a human can review instead of re-reading the transcript.
\textsc{Charter} is a single stdlib-only Python file, integrates with CI systems,
pre-commit, and agent frameworks (including a stop-hook mode that prevents an
agent from declaring completion while required checks fail), and never calls a
model in the verification path.

The empirical question is whether gating generation on such a contract actually
changes outcomes, and at what cost. No public benchmark directly measures
``deliverable governance for long-horizon runs,'' so we anchor the claim on the
closest recognized proxy with verifiable acceptance criteria: IFEval
\citep{zhou2023ifeval}, \NPrompts{} prompts each annotated with programmatically
checkable constraints (word counts, required keywords, formatting, etc.). We
compile each prompt's constraints into a \textsc{Charter} contract using only the
tool's generic check primitives, and compare single-shot generation against a
gate-and-revise loop in a paired design. Three results stand out. First, the gate
recovers a large fraction of the compliance that a capable model misses
single-shot (+\DeltaPromptStrict{} points prompt-level strict accuracy). Second,
the recovery is concentrated on constraints the contract can express: instruction
types our compiler covers exactly improve by \ExactDelta{} points, while types it
cannot express move by only \UncoveredDelta{} points---the gate's value is bounded
by its coverage, which we measure and report. Third, the cost is modest and
predictable (\MeanRevisions{} extra model calls per prompt on average, incurred
only when a check fails).

We do not claim that contract gating improves a model's underlying capabilities,
nor that IFEval is a long-horizon benchmark. The claim is narrower and, we
believe, more useful: when acceptance criteria can be written down as
deterministic checks---which covers a surprisingly large share of what ``done''
means for real deliverables---putting those checks in the loop converts silent
non-compliance into either compliance or an explicit, auditable failure.

Our contributions are: (1) a problem framing---deliverable contracts---for
governing the outputs of long-horizon agent runs, separating what is verifiable
mechanically from what requires human or model judgment; (2) \textsc{Charter}, a
production-quality, dependency-free implementation designed for adoption inside
existing CI and agent-loop infrastructure; (3) a paired evaluation on a public
benchmark quantifying both the benefit and the boundary (coverage) of
contract-gated generation, with all logs, contracts, and scoring code released.

\section{Related Work}

\paragraph{Verifiable instruction following.} IFEval~\citep{zhou2023ifeval}
established the practice of evaluating LLMs on constraints that can be checked by
a program rather than a judge, and is part of standard open evaluation suites. We
use it not as a capability benchmark but as a source of ground-truth acceptance
criteria against which a governance mechanism can be measured with no judge in
the loop.

\paragraph{Self-correction and refinement.} Self-Refine~\citep{madaan2023selfrefine}
and related work show that models can improve their own outputs given feedback;
\citet{huang2024cannot} show that \emph{intrinsic} self-correction---without
external signals---often fails or degrades performance, and \citet{pan2023survey}
survey the space of correction strategies. Our results are consistent with this
literature: the gains in our treatment arm come from cheap, \emph{external},
deterministic verification compiled from the task specification, not from asking
the model to critique itself. \textsc{Charter}'s contribution is not a new
correction algorithm but the contract abstraction and tooling that make external
verification deployable in real agent workflows.

\paragraph{Constrained and guided generation.} Grammar- and automaton-constrained
decoding~\citep{willard2023outlines} guarantees syntactic properties (e.g.,
JSON well-formedness) at decoding time, but requires logit-level access and
cannot express cross-artifact, filesystem-level criteria (``the report cites at
least five sources and the tests pass''). Contract gating operates after
generation, treats the generator as a black box, and composes with constrained
decoding where available.

\paragraph{Evaluating agents.} Agent benchmarks such as
SWE-bench~\citep{jimenez2024swebench}, GAIA~\citep{mialon2023gaia}, and
AgentBench~\citep{liu2023agentbench} score end-task success of agent systems.
They answer ``how capable is this agent?''; we address the complementary,
deployment-time question ``did \emph{this run} deliver what \emph{I} specified,
and can I verify that without reading everything?''. LLM-as-a-judge
evaluation~\citep{zheng2023judging} is the prevailing answer for criteria that
resist mechanical checks; \textsc{Charter} deliberately handles the deterministic
complement and can invoke judges through its \texttt{command} check where teams
choose to.

\paragraph{Design by contract and CI.} The contract framing borrows directly from
design-by-contract in software engineering~\citep{meyer1992dbc} and from the
practice of gating merges and releases on automated checks. Our observation is
that agent deliverables currently lack the equivalent layer, even though the
artifacts they produce (files in a workspace) are exactly what such tooling is
good at checking.

\section{Charter: deliverable contracts for agent runs}
\label{sec:design}

\paragraph{Contract model.} A charter is a TOML or JSON document with a list of
\emph{deliverables}, each with an identifier, an optional path or glob, a
required/optional flag, and a list of \emph{checks}. Checks are drawn from a
small set of primitives chosen to be deterministic, dependency-free, and fast:
\texttt{exists}, \texttt{not\_exists}, \texttt{file\_count}, \texttt{min\_bytes},
\texttt{max\_bytes}, \texttt{contains} / \texttt{not\_contains} (literal or regex,
with occurrence bounds), \texttt{word\_count}, \texttt{line\_count},
\texttt{json\_valid} (with required dotted keys), \texttt{markdown\_sections},
\texttt{no\_placeholder} (TODO/FIXME/lorem-ipsum style scaffolding), and
\texttt{command}, which delegates to any executable check the team already
trusts (test suites, linters, schema validators) and subsumes domain-specific
verification. Checks carry a severity (\texttt{error} or \texttt{warn}); a run
fails if any required deliverable has a failing error-severity check.

\paragraph{Three uses of one artifact.} The same contract is (i) \emph{rendered
for the agent}: \texttt{charter explain} produces a human- and agent-readable
statement of what must exist and how it will be judged, suitable for inclusion in
the system or task prompt; (ii) \emph{enforced}: \texttt{charter check} evaluates
the workspace and exits 0 or 1, prints a structured failure report (also
available as JSON), and has a hook mode in which a failing contract blocks an
agent's attempt to finish and feeds the failure list back into the loop; and
(iii) \emph{audited}: \texttt{charter report} emits per-check verdicts together
with SHA-256 hashes, sizes, and timestamps of the artifacts, so a reviewer can
trust the verdict without re-deriving it. Because verification is deterministic
and local, the gate can run on every agent iteration at negligible cost, and the
same contract yields the same verdict for the agent, the CI system, and the
human.

\paragraph{What contracts deliberately do not do.} Charter does not judge
qualities that require taste or domain expertise (is the analysis insightful?),
does not verify factual grounding, and does not attempt to be a sandbox or
permission system. It narrows the surface a human must review, and makes
explicit---in a reviewable file rather than in prompt history---which acceptance
criteria are mechanically enforced.

\section{Compiling a public benchmark's constraints into contracts}
\label{sec:compile}

To measure the effect of contract gating we need tasks whose acceptance criteria
are known and independently scoreable. IFEval provides \NPrompts{} prompts with
\NInstrInstances{} verifiable instruction instances across 25 types. A small
compiler (\texttt{compile\_contract.py}, 200 lines) maps each instruction onto
\textsc{Charter}'s generic primitives and labels every instance with a
\emph{coverage class}:

\begin{itemize}
  \item \textbf{exact} (\ExactN{} instances, \ExactPct\%): the compiled check is
  logically equivalent to the official checker, e.g.\ keyword presence/frequency,
  forbidden words, no-comma, postscript, title markers, end-phrase, quotation
  wrapping, placeholder counts.
  \item \textbf{approximate} (\ApproxN{} instances, \ApproxPct\%): the check
  captures the constraint with minor semantic differences, e.g.\ whitespace word
  counts versus the official tokenizer, bullet/paragraph counting via regular
  expressions, all-caps/all-lowercase via character classes, JSON validity
  without markdown-fence stripping.
  \item \textbf{uncovered} (\UncoveredN{} instances, \UncoveredPct\%): not
  expressible with the generic primitives---response language detection, sentence
  counting, and first-word-of-$n$th-paragraph constraints.
\end{itemize}

The compiler is part of the released harness. Coverage is reported per instance
in the run logs, and Section~\ref{sec:results} breaks results down by class. We
emphasize the design choice this reflects: in deployment, the contract author
decides what is worth checking mechanically; here the benchmark's own constraint
annotations play that role, which makes the overlap between gate and scorer
intentional rather than confounding---the experiment measures how much of the
specified-and-checkable requirements a strong model silently misses, and how much
of that a contract gate recovers.

\section{Experimental setup}
\label{sec:setup}

\paragraph{Arms.} \textbf{Control}: one model call per prompt; the response is
used as-is. \textbf{Treatment}: the treatment starts from the \emph{same}
control response (exact pairing, isolating the effect of the gate from sampling
noise); the compiled contract is evaluated by the real \texttt{charter check}
CLI on the response written to a file; if any check fails, the structured
failure report is appended to a revision request and the model produces a new
response, for at most two revisions; the gate keeps the attempt with the fewest
failing checks. Prompts whose initial response already passes the contract incur
no additional calls and are byte-identical across arms.

\paragraph{Model and harness.} Both arms use \texttt{claude-haiku-4-5} (Anthropic)
called through the Claude Code CLI in non-interactive print mode with a minimal
``helpful assistant'' system prompt and no tool access; the identical stack is
used for both arms, and all calls are logged verbatim. We chose a small frontier
model deliberately: it is representative of the cost tier at which long-horizon
loops actually run, and it leaves measurable headroom on IFEval.

\paragraph{Scoring and statistics.} All accuracy numbers are computed by the
official, unmodified \texttt{instruction\_following\_eval} library released with
IFEval (vendored in our repository), in both strict and loose modes, at the
prompt level (all instructions of a prompt satisfied) and the instruction level.
Because arms are paired by prompt, we report McNemar's exact test on discordant
pairs and a nonparametric bootstrap (10{,}000 resamples) for the 95\% confidence
interval of the accuracy difference. The full pipeline---harness, contracts, raw
responses for every attempt, scoring and analysis scripts---is committed alongside
this paper, and every number below is generated programmatically from those
artifacts.

\section{Results}
\label{sec:results}

\begin{table}[t]
  \caption{Paired comparison on IFEval (\NPromptsPaired{} prompts), official
  scorer. The treatment differs from control only where a compiled
  \textsc{Charter} check failed and triggered revision.}
  \label{tab:main}
  \centering
  \begin{tabular}{lccc}
    \toprule
    Metric & Control & + \textsc{Charter} gate & $\Delta$ \\
    \midrule
    Prompt-level strict accuracy (\%) & \ControlPromptStrict & \TreatmentPromptStrict & +\DeltaPromptStrict \\
    Prompt-level loose accuracy (\%) & \ControlPromptLoose & \TreatmentPromptLoose & +\DeltaPromptLoose \\
    Instruction-level strict accuracy (\%) & \ControlInstrStrict & \TreatmentInstrStrict & --- \\
    \bottomrule
  \end{tabular}
\end{table}

\begin{table}[t]
  \caption{Instruction-level strict accuracy by \textsc{Charter} coverage class.
  The gate improves what it can check; constraints it cannot express are
  essentially unchanged.}
  \label{tab:coverage}
  \centering
  \begin{tabular}{lcccc}
    \toprule
    Coverage class & Instances & Control (\%) & + gate (\%) & $\Delta$ \\
    \midrule
    Exact        & \ExactN     & \ExactControl     & \ExactTreat     & +\ExactDelta \\
    Approximate  & \ApproxN    & \ApproxControl    & \ApproxTreat    & +\ApproxDelta \\
    Uncovered    & \UncoveredN & \UncoveredControl & \UncoveredTreat & +\UncoveredDelta \\
    \midrule
    All          & \AllN       & \AllControl       & \AllTreat       & +\AllDelta \\
    \bottomrule
  \end{tabular}
\end{table}

\paragraph{Main effect.} Table~\ref{tab:main} summarizes the paired comparison.
Single-shot generation with a capable small model satisfies all of a prompt's
verifiable constraints \ControlPromptStrict\% of the time under the official
strict scorer; under the \textsc{Charter} gate this rises to
\TreatmentPromptStrict\% (+\DeltaPromptStrict{} points, 95\% bootstrap CI
[+\DeltaCILow, +\DeltaCIHigh], McNemar exact $p = \McNemarP$). The discordant
pairs are heavily one-sided: \DiscordantTreatOnly{} prompts are satisfied only
under the gate versus \DiscordantControlOnly{} only under control. Loose-mode
accuracy moves similarly (+\DeltaPromptLoose{} points), indicating the gains are
not artifacts of strict-mode formatting technicalities. Both regressions
(prompts satisfied only under control) occurred on approximately-covered
constraints, where a revision satisfied the compiled approximation but violated
the official checker's stricter semantics---a concrete cost of approximate
coverage that Section~\ref{sec:limitations} returns to.

\paragraph{The gain follows coverage.} Table~\ref{tab:coverage} breaks
instruction-level strict accuracy down by the compiler's coverage class.
Exactly-covered constraint types improve by \ExactDelta{} points and
approximately-covered types by \ApproxDelta{} points, while uncovered types move
by only \UncoveredDelta{} points---within noise of no change, as expected since
the gate cannot observe them. This is the central honest finding: contract
gating is not a general capability amplifier; it converts \emph{specified and
checkable} requirements from ``usually followed'' to ``almost always followed,''
and does nothing for requirements left outside the contract.

\paragraph{Gate dynamics and cost.} Of \NPromptsPaired{} prompts,
\GateInitialFail{} failed at least one compiled check on the initial response;
after at most two gated revisions, \GateFinalFail{} still failed at least one
check. The gate selected a revised response for \GateSelectedRevision{} prompts;
for all others the original response was already contract-clean and the arms are
identical. The loop cost \ExtraCalls{} additional model calls in total
(\MeanRevisions{} per prompt on average, concentrated on the
\PromptsWithRevision{} prompts that needed revision), and mean response length
was essentially unchanged (\MeanCharsControl{} versus \MeanCharsTreatment{}
characters), i.e.\ the gate did not induce padding. There is no
verification-time model cost: checks are regular expressions and file
inspections.

\paragraph{Qualitative observation from dogfooding.} The research run that
produced this paper was itself governed by a charter covering its required
deliverables. The gate caught, among other things, marketing copy that tripped a
placeholder-text check because it \emph{described} placeholder text---a
false positive that was resolved by narrowing the contract, illustrating the
intended workflow: acceptance criteria are tuned in a reviewable file, not in
prompt history.

\section{Limitations}
\label{sec:limitations}

\textbf{Benchmark proxy.} IFEval is single-turn instruction following, not a
long-horizon agent benchmark; we use it because it is public, recognized, and has
verifiable acceptance criteria, but the transfer of effect sizes to multi-hour,
multi-artifact runs is an extrapolation. Building a public benchmark of
contract-governed long-horizon deliverables is future work. \textbf{Gate--scorer
overlap.} The treatment optimizes checks compiled from the same constraint
annotations the scorer uses; results quantify recovery of specified-and-checkable
requirements, not generalization to unspecified ones (Table~\ref{tab:coverage}
makes the boundary explicit). \textbf{Coverage ceiling.} \UncoveredPct\% of
constraint instances were not expressible with generic primitives, and
approximate checks can both miss violations and raise false alarms; richer
primitives (or \texttt{command}-delegated domain validators) shift but do not
remove this boundary. \textbf{Single model, single revision policy.} We evaluate
one model at one cost tier with at most two revisions; gains may differ for
stronger models (less headroom) or weaker ones (revisions may fail more often).
\textbf{Contracts can be gamed.} A model optimizing against a literal check can
satisfy it in unintended ways (e.g., padding to reach a word count); deterministic
gates should be combined with human or model review for adversarial settings.

\section{Conclusion}

Treating the deliverables of long-horizon agent runs as the subject of a small,
declarative, machine-checkable contract is both practically deployable today and
measurably effective at eliminating the silent non-compliance that currently
forces humans to re-read everything an agent produces. The boundary of the
approach is exactly its coverage---a property we can measure, report, and expand
deliberately---which we argue is the right shape for a governance mechanism: it
is auditable about what it does and does not guarantee. \textsc{Charter}, the
benchmark harness, raw logs, and this paper's numbers are released for
reproduction and adoption.

\begin{thebibliography}{12}

\bibitem[Huang et~al.(2024)]{huang2024cannot}
Jie Huang, Xinyun Chen, Swaroop Mishra, Huaixiu~Steven Zheng, Adams~Wei Yu,
Xinying Song, and Denny Zhou.
\newblock Large language models cannot self-correct reasoning yet.
\newblock In \emph{International Conference on Learning Representations (ICLR)}, 2024.
\newblock arXiv:2310.01798.

\bibitem[Jimenez et~al.(2024)]{jimenez2024swebench}
Carlos~E. Jimenez, John Yang, Alexander Wettig, Shunyu Yao, Kexin Pei, Ofir
Press, and Karthik Narasimhan.
\newblock {SWE}-bench: Can language models resolve real-world {GitHub} issues?
\newblock In \emph{International Conference on Learning Representations (ICLR)}, 2024.
\newblock arXiv:2310.06770.

\bibitem[Liu et~al.(2023)]{liu2023agentbench}
Xiao Liu, Hao Yu, Hanchen Zhang, Yifan Xu, Xuanyu Lei, Hanyu Lai, Yu Gu, Hangliang
Ding, Kaiwen Men, Kejuan Yang, et~al.
\newblock Agentbench: Evaluating {LLMs} as agents.
\newblock \emph{arXiv preprint arXiv:2308.03688}, 2023.

\bibitem[Madaan et~al.(2023)]{madaan2023selfrefine}
Aman Madaan, Niket Tandon, Prakhar Gupta, Skyler Hallinan, Luyu Gao, Sarah
Wiegreffe, Uri Alon, Nouha Dziri, Shrimai Prabhumoye, Yiming Yang, et~al.
\newblock Self-refine: Iterative refinement with self-feedback.
\newblock In \emph{Advances in Neural Information Processing Systems (NeurIPS)}, 2023.
\newblock arXiv:2303.17651.

\bibitem[Meyer(1992)]{meyer1992dbc}
Bertrand Meyer.
\newblock Applying ``design by contract''.
\newblock \emph{IEEE Computer}, 25(10):40--51, 1992.

\bibitem[Mialon et~al.(2023)]{mialon2023gaia}
Gr\'egoire Mialon, Cl\'ementine Fourrier, Craig Swift, Thomas Wolf, Yann LeCun,
and Thomas Scialom.
\newblock {GAIA}: a benchmark for general {AI} assistants.
\newblock \emph{arXiv preprint arXiv:2311.12983}, 2023.

\bibitem[Pan et~al.(2024)]{pan2023survey}
Liangming Pan, Michael Saxon, Wenda Xu, Deepak Nathani, Xinyi Wang, and
William~Yang Wang.
\newblock Automatically correcting large language models: Surveying the landscape
of diverse automated correction strategies.
\newblock \emph{Transactions of the Association for Computational Linguistics},
12:484--506, 2024.
\newblock arXiv:2308.03188.

\bibitem[Willard and Louf(2023)]{willard2023outlines}
Brandon~T. Willard and R\'emi Louf.
\newblock Efficient guided generation for large language models.
\newblock \emph{arXiv preprint arXiv:2307.09702}, 2023.

\bibitem[Zheng et~al.(2023)]{zheng2023judging}
Lianmin Zheng, Wei-Lin Chiang, Ying Sheng, Siyuan Zhuang, Zhanghao Wu, Yonghao
Zhuang, Zi Lin, Zhuohan Li, Dacheng Li, Eric Xing, et~al.
\newblock Judging {LLM}-as-a-judge with {MT-Bench} and {Chatbot Arena}.
\newblock In \emph{Advances in Neural Information Processing Systems (NeurIPS)}, 2023.
\newblock arXiv:2306.05685.

\bibitem[Zhou et~al.(2023)]{zhou2023ifeval}
Jeffrey Zhou, Tianjian Lu, Swaroop Mishra, Siddhartha Brahma, Sujoy Basu, Yi Luan,
Denny Zhou, and Le Hou.
\newblock Instruction-following evaluation for large language models.
\newblock \emph{arXiv preprint arXiv:2311.07911}, 2023.

\end{thebibliography}

\end{document}
